\begin{table}[!htbp]
\centering
\caption{Where the gap arises: applications versus decisions}
\label{tab:pipeline}
\footnotesize
\begin{adjustbox}{max width=\linewidth}
\begin{tabular}{@{}lccccccc@{}}
\toprule
 & \multicolumn{3}{c}{Application margin} & \multicolumn{3}{c}{Share of decided subapplications} & Award \\
\cmidrule(lr){2-4}\cmidrule(lr){5-7}\cmidrule(lr){8-8}
 & Applied & \shortstack{Log(1+\\number)} & \shortstack{County gov.\\applied} & Approved & Withdrawn & Denied & \shortstack{County gov.\\funded} \\
 & (1) & (2) & (3) & (4) & (5) & (6) & (7) \\
\midrule
Small county & $-5.37^{***}$ & $-4.68^{***}$ & $-2.35^{***}$ & $-0.52$ & $-0.87$ & $1.40^{**}$ & $-2.38^{***}$ \\
 & $(0.98)$ & $(1.02)$ & $(0.73)$ & $(1.90)$ & $(1.82)$ & $(0.65)$ & $(0.74)$ \\
Mean, larger counties & 26.1 & 25.9 & 11.1 & 86.6 & 11.7 & 1.7 & 11.4 \\
Observations & 12,038 & 12,038 & 12,038 & 2,515 & 2,515 & 2,515 & 12,038 \\
Counties (clusters) & 2,762 & 2,762 & 2,762 & 1,459 & 1,459 & 1,459 & 2,762 \\
\bottomrule
\end{tabular}
\end{adjustbox}
\begin{minipage}{0.95\linewidth}\vspace{4pt}\scriptsize
\textit{Note:} All columns use the column (5) specification of Table~\ref{tab:baseline} (declaration fixed effects; hazard, demographic, and local-government-count controls). Coefficients multiplied by 100. Columns (1)--(3) and (7) use all county--declarations; ``county gov.'' restricts to subapplications or projects whose applicant is the county government itself, of which every county has exactly one. Columns (4)--(6) restrict to county--declarations with at least one decided subapplication and measure the share approved, withdrawn, or denied. Standard errors clustered by county. At the subapplication level (same controls), small-county subapplications request 0.102 log points less funding (SE 0.094) and report benefit--cost ratios 0.069 log points lower (SE 0.063). $^{*}p<0.1$; $^{**}p<0.05$; $^{***}p<0.01$.
\end{minipage}
\end{table}
